% Industry resume, Grover Aruquipa
% Compile with pdfLaTeX (also works on Overleaf). One column on purpose:
% applicant tracking systems read single-column text far more reliably than sidebars.
\documentclass[10pt,letterpaper]{article}
\usepackage[T1]{fontenc}
\usepackage[utf8]{inputenc}
\usepackage[scaled=0.95]{helvet}
\renewcommand{\familydefault}{\sfdefault}
\usepackage[top=1.0cm,bottom=0.8cm,left=1.5cm,right=1.5cm]{geometry}
\usepackage{enumitem}
\usepackage{titlesec}
\usepackage{xcolor}
\usepackage{microtype}
\usepackage[hidelinks]{hyperref}

\definecolor{ink}{HTML}{111111}
\definecolor{accent}{HTML}{0D2B61}
\definecolor{muted}{HTML}{555555}
\color{ink}
\pagestyle{empty}
\setlength{\parindent}{0pt}
\setlength{\parskip}{0pt}

\titleformat{\section}{\normalsize\bfseries\color{accent}}{}{0em}{\MakeUppercase}[\vspace{-6pt}{\color{accent}\rule{\linewidth}{0.6pt}}]
\titlespacing*{\section}{0pt}{7pt}{3pt}
\setlist[itemize]{leftmargin=1.1em,itemsep=0.5pt,topsep=1pt,parsep=0pt,label=\textbullet}

% \role{Title}{Organization}{Location}{Dates}
\newcommand{\role}[4]{\textbf{#1}, #2 \hfill {\color{muted}#3 \textbar\ #4}\par}

\begin{document}

%================ HEADER ================
{\LARGE\bfseries Grover Aruquipa}\par\vspace{2pt}
{\large Robotics Software Developer \textbar\ ROS~2, MoveIt~2, Simulation and Perception in C++ and Python}\par\vspace{3pt}
{\small Qu\'ebec, QC, Canada \ \textbullet\ \
\href{mailto:grover.aruquipa.9@gmail.com}{grover.aruquipa.9@gmail.com} \ \textbullet\ \
\href{https://www.linkedin.com/in/grover-aruquipa-996917126}{linkedin.com/in/grover-aruquipa-996917126} \ \textbullet\ \
\href{https://groveraruquipa.github.io}{groveraruquipa.github.io}}\par


%================ SUMMARY ================
\section{Summary}
Robotics engineer finishing a PhD (Universit\'e Laval, defense early 2027) who writes the software that runs real robots:
kinematics, motion and collision checking, simulation, real-time control and perception. I have brought three robots from
model to working hardware, written their control in C++ and MATLAB at 500\,Hz, built a ROS~2 and MoveIt~2 stack for a
9-DoF robot, and simulated them in MuJoCo and Isaac Sim. Previously computer vision engineer and data science intern in
industry. Eligible for an open Post-Graduation Work Permit in Canada.

%================ SKILLS ================
\section{Technical Skills}
\begin{itemize}
\item \textbf{Robotics software:} ROS and ROS~2, MoveIt and MoveIt~2, ros2\_control; motion planning and collision checking;
kinematics, Jacobians and singularity avoidance; Pinocchio.
\item \textbf{Simulation:} NVIDIA Isaac Sim, MuJoCo, Simscape and Simulink; simulation scenes and robot models (URDF) built
from my own CAD.
\item \textbf{Programming:} C++ (real time), Python, MATLAB; Git, Linux, Docker devcontainers (GitHub Codespaces).
\item \textbf{Perception and ML:} OpenCV, object detection in industrial scenes, PyTorch, TensorFlow; machine learning on
sensor data (vibration, high-voltage equipment).
\item \textbf{Hardware integration:} real-time multi-axis control at 500\,Hz (OPAL-RT, Simulink Real-Time), motors, drives
and encoders; SolidWorks; debugging software and hardware together on the bench.
\end{itemize}

%================ EXPERIENCE ================
\section{Experience}
\role{Doctoral Researcher}{Robotics Laboratory, Universit\'e Laval}{Qu\'ebec, Canada}{2023 -- present}
\begin{itemize}
\item Wrote the kinematics, redundancy resolution and real-time control software for three robots with 4, 7 and 8
actuators, and debugged them on the hardware until they met their targets.
\item \textbf{ROS~2 and MoveIt~2:} built a ROS~2 package for a 9-DoF parallel robot with grasping (robot description,
ros2\_control, MoveIt~2 planning) in a Docker devcontainer.
\item \textbf{Simulation:} simulated my robots in MuJoCo to test motions and grasps before running them on hardware, and
set up Isaac Sim on a cloud GPU machine to experiment with robot models.
\item \textbf{Collision checking:} developed an analytical method that keeps unlimited rotations free of collisions and
singularities, used to plan the motions of the robots (under review, Mechanism and Machine Theory).
\item \textbf{Results:} coaxial wrist-gripper turned its tool ten full revolutions with sub-degree tracking at 500\,Hz;
Delta robot in progress that estimates grasp force from motor currents. Day-to-day lead of two graduate students.
\end{itemize}

\vspace{3pt}
\role{Visiting Researcher}{Institute of Robotics, Johannes Kepler University Linz}{Linz, Austria}{Apr. -- May 2026}
\begin{itemize}
\item Developed singularity and grasp-quality indices for manipulators with integrated grasping, hosted by Univ.-Prof.~A.~M\"uller (Mitacs Globalink Research Award).
\end{itemize}

\vspace{3pt}
\role{Research Intern}{FEMTO-ST (CNRS)}{Besan\c{c}on, France}{2022 -- 2023}
\begin{itemize}
\item Designed and modeled a redundant dexterous parallel robot with two unlimited rotations for micro-manipulation, and wrote its real-time C++ control software; clean-room fabrication projects.
\end{itemize}

\vspace{3pt}
\role{Computer Vision Engineer}{Sencon Srl}{La Paz, Bolivia}{2021}
{\small Industrial computer vision systems, specializing in scaffolding detection in industrial environments.}\par\vspace{2pt}
\role{Data Science Intern}{Hitachi ABB Power Grids}{Zurich, Switzerland}{2020}
{\small Machine learning for the analysis of high-voltage circuit breakers. Also data scientist for KTH (Sweden), 2021.}

%================ EDUCATION ================
\section{Education}
\textbf{PhD, Mechanical Engineering}, Universit\'e Laval \hfill {\color{muted}Qu\'ebec, Canada \textbar\ 2023 -- 2027 (expected)}\par
{\small Kinematically redundant parallel robots with unlimited rotations and integrated grasping. Advisors: C.~Gosselin, P.~Lambert.}\par\vspace{2pt}
\textbf{MSc, Smart Integrated Systems}, Universit\'e Bourgogne--Franche-Comt\'e \hfill {\color{muted}Besan\c{c}on, France \textbar\ 2021 -- 2023}\par\vspace{2pt}
\textbf{BEng, Mechatronics Engineering}, Bolivian Catholic University \hfill {\color{muted}La Paz, Bolivia \textbar\ 2014 -- 2019}\par
{\small Thesis: reconfigurable Delta robot for product picking.}

%================ AWARDS ================
\section{Awards and Languages}
\begin{itemize}
\item 1st Place, Graduate Division, ASME Student Mechanisms and Robotics Design Competition (2025), team lead.
\item Finalist, MassRobotics competition, Boston, with one of my robot prototypes, team lead.
\item Mitacs Globalink Research Award (2026); 1st Place, Student Paper Competition, Argentine Robotics Conference (2019).
\item Languages: Spanish (native), English (full professional), French (intermediate).
\end{itemize}

\newpage
%================ PUBLICATIONS
\section{Selected Publications}
{\small First author on all papers. Full list at groveraruquipa.github.io.}\par\vspace{2pt}
\begin{itemize}
\item \textbf{G.~Aruquipa}, P.~Lambert, C.~Gosselin, ``Analysis and design of a novel (6+1)-DoF parallel robot with unlimited rotation and grasping capabilities,'' \emph{Mechanism and Machine Theory}, vol.~224, art.~106438, 2026.
\item \textbf{G.~Aruquipa}, P.~Lambert, C.~Gosselin, ``Kinematic analysis and design of a novel 9-DoF parallel robot with grasping capabilities,'' \emph{CCToMM Symposium on Mechanisms, Machines and Mechatronics}, 2025, pp.~117--127.
\item \textbf{G.~Aruquipa}, P.~Lambert, C.~Gosselin, ``An analytical framework for collision-free and singularity-free unlimited rotation in spatial parallel mechanisms,'' under review, \emph{Mechanism and Machine Theory}.
\item \textbf{G.~Aruquipa}, V.~Muralidharan, P.~Lambert, C.~Gosselin, ``Analysis and design of a novel compact four-degree-of-freedom parallel robot with unlimited rotation,'' \emph{ASME Journal of Mechanisms and Robotics}, vol.~18, no.~6, art.~061001, 2026.
\item \textbf{G.~Aruquipa}, F.~Diaz, ``An IoT architecture based on the control of a bio-inspired manufacturing system for the detection of anomalies with vibration sensors,'' \emph{Procedia Computer Science}, vol.~200, pp.~438--450, 2022.
\end{itemize}

%================ REFERENCES ================
\section{References}
\begin{itemize}
\item \textbf{Prof.~Cl\'ement Gosselin}, Full Professor, Department of Mechanical Engineering, Universit\'e Laval, Qu\'ebec,
Canada. PhD supervisor. \href{mailto:clement.gosselin@gmc.ulaval.ca}{clement.gosselin@gmc.ulaval.ca}
\item \textbf{Prof.~Patrice Lambert}, Assistant Professor, Department of Mechanical Engineering, Universit\'e Laval,
Qu\'ebec, Canada. PhD co-supervisor. \href{mailto:patrice.lambert@gmc.ulaval.ca}{patrice.lambert@gmc.ulaval.ca}
\item \textbf{Univ.-Prof.~Dr.-Ing.~Andreas M\"uller}, Head of the Institute of Robotics, Johannes Kepler University Linz,
Austria. Host during my research stay (2026). \href{mailto:a.mueller@jku.at}{a.mueller@jku.at}
\end{itemize}

\end{document}
